\chapter{Introdução}
\label{sec:introducao}

O rádio definido por software (\textit{Software-Defined Radio}, SDR) é uma arquitetura de sistema de rádio em que funções tradicionalmente resolvidas por hardware dedicado — sintonia, filtragem e demodulação — são implementadas em software. O sinal é digitalizado o mais cedo possível na cadeia de recepção, por um conversor analógico-digital, o que torna o sistema reconfigurável: mudar o canal sintonizado, a largura de banda ou o esquema de demodulação passa a ser uma questão de trocar coeficientes e rotinas, e não de substituir circuitos.

Como consequência dessa escolha de arquitetura, uma fatia relativamente larga do espectro é digitalizada de uma só vez, e a separação entre os diversos canais presentes nessa faixa deixa de ser um problema de circuitos e passa a ser um problema puramente algorítmico — resolvido por filtros digitais.

\section{Motivação}

Considera-se um receptor que digitaliza uma banda de $F_s = 1{,}92$~MHz contendo múltiplos canais adjacentes. Nessa janela há vários canais transmitindo lado a lado, como estações de rádio vizinhas. Após a conversão descendente digital (\textit{Digital Down-Conversion}, DDC), o canal de interesse é trazido para a banda base, resultando em um sinal complexo em fase e quadratura (I/Q) centrado em $0$~Hz e ocupando aproximadamente $\pm 100$~kHz. O DDC, porém, apenas escolhe um dos canais: os vizinhos continuam presentes no sinal digitalizado. Um canal adjacente permanece centrado em $200$~kHz e um interferente de amplitude elevada permanece centrado em $450$~kHz, ambos dentro da faixa digitalizada.

Como o conteúdo útil do sinal ocupa uma banda muito menor do que a taxa de amostragem original, o sinal encontra-se significativamente sobreamostrado. Isso representa um custo computacional desnecessário para o bloco subsequente da cadeia de recepção — tipicamente um classificador automático de modulação, que opera sobre a constelação I/Q do sinal recebido. Propõe-se, portanto, reduzir a taxa de amostragem por meio de uma decimação de fator $M = 4$, levando a taxa efetiva de $1{,}92$~MHz para $480$~kHz.

Antes dessa redução, entretanto, um único filtro digital passa-baixas precisa cumprir duas funções simultâneas: (i) a \textbf{seleção de canal}, isolando o sinal útil (de $0$ a $100$~kHz) e rejeitando o canal adjacente a partir de $200$~kHz; e (ii) o \textbf{anti-\textit{aliasing}} para a decimação, eliminando toda energia que, ao reduzir a taxa, dobraria sobre o canal útil — em particular, o interferente em $450$~kHz. Trata-se de um filtro passa-baixas no sentido clássico: deixa passar as frequências baixas correspondentes ao canal em banda base e atenua as altas, correspondentes aos canais vizinhos e ao interferente.

\section{Justificativa do fator de decimação}

A escolha de $M = 4$ não é arbitrária: decorre de um compromisso entre o limite imposto pelo teorema da amostragem, o custo computacional do filtro e a margem operacional exigida pelo bloco seguinte. Como o sinal em banda base é complexo, a taxa mínima que preserva a informação corresponde à própria largura de banda ocupada, e não ao seu dobro; um fator $M = 8$ ($F_s' = 240$~kHz) ainda seria teoricamente admissível, pois o canal caberia dentro da nova Nyquist de $\pm 120$~kHz — deixando, porém, apenas $20$~kHz de banda de guarda de cada lado.

O custo aparece na condição de anti-\textit{aliasing}: para proteger o canal de $\pm 100$~kHz, o filtro precisa atenuar toda energia nas faixas $k\!\cdot\!(F_s/M) \pm 100$~kHz, para $k$ inteiro. Com $M = 4$, a primeira delas ($380$ a $580$~kHz) já está inteiramente coberta pela banda de rejeição exigida para bloquear o canal adjacente, e o anti-\textit{aliasing} é obtido sem exigência adicional sobre o filtro. Com $M = 8$, ela desce para o intervalo de $140$ a $340$~kHz, obrigando a banda de rejeição a começar em $140$~kHz e estreitando a faixa de transição de $100$ para $40$~kHz — o que, sendo a ordem de um FIR aproximadamente inversa a essa largura, elevaria o comprimento do filtro dos $93$ coeficientes do projeto para algo próximo de $230$. A economia na taxa de saída seria paga na filtragem.

Soma-se a isso a margem operacional: o classificador de modulação opera sobre a constelação I/Q e requer sobreamostragem em relação à taxa de símbolo — para recuperação de sincronismo, filtro casado e interpolação do instante ótimo de decisão —, e desvios residuais de frequência do estágio de DDC, bem como efeito Doppler, deslocam ligeiramente a posição espectral do canal. Os $20$~kHz de guarda de $M = 8$ seriam frágeis frente a esses desvios, enquanto $M = 4$ entrega um fator de sobreamostragem de $2{,}4$ vezes sobre a banda ocupada. O fator de decimação é, portanto, uma decisão de sistema, e não um número arbitrário do enunciado.